% ===========================================================================
%  四元数姿态误差
% ===========================================================================

\section{四元数姿态误差与大姿态控制}
\label{sec:quat_control}

\subsection{误差定义必须与乘法方向一致}

四元数避免 3--2--1 欧拉角在 $\theta=\pm90^\circ$ 处的参数奇异，适合垂直姿态和
大角度指令\cite{kuipers1999quaternions,markley2014fundamentals,fresk2013full}。
但“四元数控制”并不自动消除符号错误；必须同时规定四元数表示的旋转方向、乘法顺序和
角速度所在坐标系。

固件采用标量在前的单位四元数，并定义
\begin{equation}
\bm q_e=\bm q^{-1}\otimes\bm q_d
=\begin{bmatrix}\eta_e&\bm\epsilon_e^T\end{bmatrix}^{T}.
\label{eq:quat_err}
\end{equation}
该误差表示从当前姿态到目标姿态的短弧旋转，并与固件中的正反馈符号组合使用。
Web 垂直模式当前使用相反乘法顺序
$\bm q_d^{-1}\otimes\bm q$，随后在外环加入负号。两种约定在各自闭环中都可能成立，
但它们的矢量部分通常互为相反方向且表达坐标不同，不能从一个实现复制误差公式、
再从另一个实现复制控制增益。

\subsection{短弧旋转向量}

单位四元数具有双覆盖：$\bm q_e$ 与 $-\bm q_e$ 表示同一姿态。固件取
$s_e=\operatorname{sgn}(\eta_e)$，并由
\begin{equation}
\bm e_R=s_e
\frac{2\operatorname{atan2}(\|\bm\epsilon_e\|,|\eta_e|)}
{\|\bm\epsilon_e\|}\bm\epsilon_e
\label{eq:quat_rotation_vector}
\end{equation}
提取短弧旋转向量；小误差时 $\bm e_R\approx2s_e\bm\epsilon_e$。
使用全向量模长而不是逐轴模长，能够在多轴误差同时存在时保持统一的轴角尺度。
固件的分段阈值和 deg/rad 换算见式\eqref{eq:firmware_rotation_vector}。

\subsection{级联外环}

对静止目标或缓慢变化目标，常用外环形式为
\begin{equation}
\bm\omega_{\mathrm{ref}}
=\bm K_R\bm e_R,
\label{eq:quat_pd}
\end{equation}
并由角速度内环跟踪。若目标姿态随时间变化，还应把参考角速度按一致坐标变换后前馈，
而不能把定点调节公式直接用于快速轨迹。固件滚转/俯仰使用式\eqref{eq:quat_pd}的比例形式，
偏航回中则使用包络航向误差，不使用 $q_{e,z}$ 的同构外环。

短弧选择会在 $180^\circ$ 误差附近引入切换面，而 $\mathrm{SO}(3)$ 的拓扑也阻止连续静态反馈
在所有姿态上实现单一全局渐近稳定平衡。因此，本文只把该控制律称为避免欧拉角奇异的
短弧局部调节器，不声称已经证明全局渐近稳定\cite{wie1989quaternion}。

\subsection{参考轨迹插值}

若两个单位四元数 $\bm q_0$、$\bm q_1$ 的内积为
$c=\bm q_0^T\bm q_1$，先在 $c<0$ 时取 $\bm q_1\leftarrow-\bm q_1$，
再用球面线性插值
\begin{equation}
\operatorname{slerp}(\bm q_0,\bm q_1;s)=
\frac{\sin[(1-s)\Omega]}{\sin\Omega}\bm q_0+
\frac{\sin(s\Omega)}{\sin\Omega}\bm q_1,
\qquad \Omega=\arccos c,
\label{eq:slerp}
\end{equation}
可生成角速度较平滑的姿态参考。该公式只定义几何路径；执行器能否跟踪还取决于
路径时标、角速度/角加速度限制、推力裕度和控制分配。早期脚本把 slerp 跟踪成功等同于
转捩能力，超出了当前证据范围。

\subsection{数值与实现检查}

四元数积分采用\cref{eq:quaternion_kinematics,eq:rk4}，每个 RK4 中间阶段先投影到单位球面，
并在最终组合后再次归一化。验证四元数控制时至少应检查单位范数、误差乘法顺序、
短弧符号、角速度坐标、目标速率前馈和 $180^\circ$ 邻域的切换行为。
当前 Node 测试能够锁定 Web 约定，固件宿主机测试能够锁定固件约定；二者尚未由同一组
端到端轨迹测试证明等价。

\subsection{工程适用域}
\label{sec:quat_sas_engineering}

误差四元数解决的是姿态表示问题，不解决执行器权限、控制分配奇异、模型不确定性或
积分饱和。对本构型，垂直姿态下仍必须同时验证推重比、$\pm15^\circ$ 固件摆角、
差速低油门退化和电机 $0.28\,\mathrm{s}$ 动态。故四元数外环是必要的表示层改进，
不是全包线控制可行性的充分条件。
